% ============================================================
% CAPITOLO D - LA SCRITTURA A QUATTRO PARTI E LA CONDOTTA DELLE VOCI
% (Versione revisionata con integrazione di errori melodici e
%  risoluzione della sensibile da fonti esterne)
% ============================================================
\chapter{La scrittura a quattro parti e la condotta delle voci}
\label{chap:condotta_voci}

Nei capitoli precedenti abbiamo analizzato gli accordi come entità verticali, statiche. La vera arte dell'armonia, tuttavia, risiede nel collegarli tra loro in modo fluido e logico. Questo processo, noto come \textbf{condotta delle voci} (o \textit{voice leading}), trasforma una sequenza di accordi statici in un discorso musicale dinamico e coerente. La scrittura a quattro parti in stile corale (Soprano, Contralto, Tenore, Basso) è il modello di studio per eccellenza per apprendere e padroneggiare questi principi.

\section{Disposizione delle voci}
\label{sec:disposizione_voci}

La scrittura a quattro parti simula un coro misto. A ogni voce è assegnata un'estensione tipica, che è bene rispettare per garantire una buona sonorità e cantabilità.

\begin{definizione}[Estensione delle voci corali (SATB)]
	\label{def:estensione_voci}
	Le estensioni standard approssimative per le quattro voci sono:
	\begin{itemize}
		\item \textbf{Soprano:} Dal Do centrale (Do') al La acuto (La'').
		\item \textbf{Contralto:} Dal Sol sotto il rigo (Sol) al Mi sul quarto spazio (Mi'').
		\item \textbf{Tenore:} Dal Do nel secondo spazio in chiave di basso (Do) al La sopra il rigo (La').
		\item \textbf{Basso:} Dal Fa sulla prima linea in chiave di basso (Fa,) al Mi sopra il rigo (Mi).
	\end{itemize}
\end{definizione}

Oltre all'estensione, è fondamentale considerare la \textbf{disposizione} delle tre voci superiori (Contralto, Tenore, Soprano). La loro distanza reciproca determina la ``geometria'' dell'accordo.

\begin{definizione}[Posizione Stretta vs. Posizione Lata]
	\label{def:posizione_stretta_lata}
	\begin{itemize}
		\item \textbf{Posizione Stretta:} La distanza tra il Soprano e il Tenore è inferiore a un'ottava. Le tre voci superiori sono il più possibile ravvicinate.
		\item \textbf{Posizione Lata (o Larga):} La distanza tra il Soprano e il Tenore è uguale o superiore a un'ottava. Le voci sono più distanziate.
	\end{itemize}
\end{definizione}

\begin{regola}[Disequazioni spaziali tra le voci]
	\label{reg:distanza_voci}
	Per preservare la coesione del tessuto corale, la distanza tra voci adiacenti non deve essere eccessiva.
	\begin{itemize}
		\item Distanza tra Soprano e Contralto $\le$ \textbf{un'ottava}.
		\item Distanza tra Contralto e Tenore $\le$ \textbf{un'ottava}.
		\item La distanza tra Tenore e Basso è più flessibile, ma si consiglia di non superare la dodicesima.
	\end{itemize}
\end{regola}

\section{Le posizioni delle triadi}
\label{sec:posizioni_triadi}

Oltre alla disposizione (stretta o lata), che riguarda la distanza tra le voci, esiste un'altra classificazione fondamentale nella scrittura a quattro parti: la \textbf{posizione} della triade.

\begin{definizione}[Posizione della Triade]
	\label{def:posizione_triade}
	La \textbf{posizione} di una triade è determinata dalla nota dell'accordo che si trova al Soprano (voce più acuta):
	\begin{itemize}
		\item \textbf{Posizione di Ottava (o Fondamentale):} La fondamentale dell'accordo è al Soprano.
		\item \textbf{Posizione di Terza:} La terza dell'accordo è al Soprano.
		\item \textbf{Posizione di Quinta:} La quinta dell'accordo è al Soprano.
	\end{itemize}
\end{definizione}

È fondamentale non confondere le \textit{posizioni} con i \textit{rivolti} (Capitolo \ref{chap:triadi}): i rivolti dipendono dalla nota al \textit{basso}, mentre le posizioni dipendono dalla nota al \textit{soprano}.

\section{Tipologie di moto}
\label{sec:tipologie_moto}

Quando due voci si muovono da un accordo al successivo, creano dei ``vettori melodici'' la cui relazione reciproca definisce il tipo di moto.

\begin{definizione}[Tipi di moto tra le parti]
	\label{def:tipi_moto}
	Esistono tre tipi di moto tra una qualsiasi coppia di voci:
	\begin{itemize}
		\item \textbf{Moto Retto (o Parallelo/Diretto):} Le due voci si muovono nella stessa direzione (entrambe salgono o entrambe scendono). Se mantengono lo stesso intervallo, si parla di \textit{moto parallelo}.
		\item \textbf{Moto Contrario:} Le due voci si muovono in direzioni opposte (una sale e l'altra scende).
		\item \textbf{Moto Obliquo:} Una voce sta ferma, mentre l'altra si muove.
	\end{itemize}
\end{definizione}

Una buona scrittura armonica mostra una netta preferenza per il \textbf{moto contrario}: esso garantisce la massima indipendenza e chiarezza melodica tra le parti.

\section{Il collegamento armonico}
\label{sec:collegamento_armonico}

Il principio guida per un collegamento armonico efficace è la \textbf{legge del minimo sforzo}: le voci devono muoversi il meno possibile per raggiungere l'accordo successivo.

\begin{regola}[Principio di minima azione]
	\label{reg:minima_azione}
	Per collegare due accordi, si seguono due principi gerarchici:
	\begin{enumerate}
		\item \textbf{Note comuni:} Se i due accordi hanno una o più note in comune, queste note dovrebbero, se possibile, rimanere ferme nella stessa voce (moto obliquo).
		\item \textbf{Grado congiunto:} Le altre voci si muovono per il percorso più breve possibile, preferibilmente per grado congiunto (intervalli di seconda).
	\end{enumerate}
\end{regola}

\subsection{Tipologie di collegamento in base alla distanza}

Il collegamento tra due triadi varia in base alla distanza tra le loro fondamentali, che determina il numero di note comuni.

\begin{itemize}
	\item \textbf{Relazione di quinta (o quarta):} Gli accordi hanno \textbf{una nota in comune}. (Esempio: I-IV, I-V, V-I).
	\item \textbf{Relazione di terza (o sesta):} Gli accordi hanno \textbf{due note in comune}. È il collegamento più fluido. (Esempio: I-VI, I-III).
	\item \textbf{Relazione di seconda (o settima):} Gli accordi \textbf{non hanno note in comune}. Richiede moto contrario obbligatorio per evitare parallelismi. (Esempio: IV-V, V-VI).
\end{itemize}

% ============================================================
% SEZIONE NUOVA: ERRORI MELODICI
% (Integrazione da armoniainrete.altervista.org/armonia-lezione-08)
% ============================================================
\section{Errori melodici}
\label{sec:errori_melodici}

Prima di affrontare gli errori armonici, è necessario considerare i vincoli che riguardano il \textbf{profilo melodico di ciascuna voce presa singolarmente}. Una linea vocale deve rispettare alcuni limiti di natura musicale ed esecutiva.

\begin{regola}[Limiti degli intervalli melodici]
	\label{reg:limiti_melodici}
	Nella condotta di una singola voce valgono le seguenti norme:
	\begin{enumerate}
		\item Sono \textbf{proibiti} intervalli melodici superiori alla \textbf{sesta minore} (ad eccezione dell'ottava giusta).
		\item Sono \textbf{proibiti} tutti gli \textbf{intervalli eccedenti e diminuiti}, con le eccezioni indicate di seguito.
		\item La \textbf{settima minore ascendente} è ammessa in due casi speciali: quando sostituisce il movimento di tono discendente tra il V e il IV grado, oppure in funzione modulante.
	\end{enumerate}
\end{regola}

\begin{osservazione}[Eccezione: seconda aumentata nel minore]
	\label{oss:seconda_aumentata}
	La \textbf{seconda aumentata} che si forma tra il VI e il VII grado della scala minore armonica è \textit{tollerata} (se impossibile da evitare), ma va evitata quando possibile. Essa nasce necessariamente dalla presenza simultanea della sopradominante naturale e della sensibile alzata.
\end{osservazione}

\begin{osservazione}[Eccezione: quarta e quinta diminuita]
	\label{oss:quarta_quinta_dim}
	Sono ammessi gli intervalli di \textbf{quarta e quinta diminuita} a condizione che il movimento sia \textbf{discendente} e che la nota successiva risalga di \textbf{semitono diatonico}. Questo schema è tipico del VII grado che funge da appoggiatura.
\end{osservazione}

\begin{errore}[Quarta eccedente melodica]
	\label{err:quarta_eccedente}
	È vietato che una voce si muova per salto di \textbf{quarta eccedente} (tritono). Questo intervallo, storicamente chiamato \textit{diabolus in musica}, risulta difficile da cantare e crea una forte discontinuità melodica.
\end{errore}

\begin{regola}[Eccezione: quarta eccedente con risoluzione]
	\label{reg:quarta_eccedente_eccezione}
	Il salto di quarta eccedente è ammesso quando la nota più acuta \textbf{o} più grave del salto si muove successivamente di \textbf{semitono diatonico} nella direzione opposta, ``risolvendo'' così la tensione del tritono melodico.
\end{regola}

% ============================================================
% FALSA RELAZIONE
% ============================================================
\subsection{La falsa relazione}
\label{subsec:falsa_relazione}

Un errore melodico di natura \textit{orizzontale}, ma che si manifesta tra due voci \textit{diverse}, è la falsa relazione.

\begin{definizione}[Falsa relazione]
	\label{def:falsa_relazione}
	Si definisce \textbf{falsa relazione} una successione armonica in cui la stessa nota appare alterata in una voce e non alterata in un'altra, in accordi consecutivi. Essa distrugge la coerenza cromatica del tessuto.
	\begin{itemize}
		\item \textbf{Falsa relazione melodica:} si verifica quando un suono viene alterato in una voce diversa da quella che l'ha cantato non alterato nell'accordo precedente.
		\item \textbf{Falsa relazione di tritono:} si ha quando due voci consecutive cantano le due note del tritono diatonico (IV e VII grado della scala) in modo tale da creare una relazione armonica ambigua tra un accordo e il successivo.
	\end{itemize}
\end{definizione}

\begin{regola}[Come evitare la falsa relazione]
	\label{reg:evitare_falsa_relazione}
	Se una nota viene alterata cromaticamente, la stessa voce che l'ha cantata non alterata deve eseguire anche la versione alterata nell'accordo successivo. Non si deve ``passare'' l'alterazione a un'altra voce.
\end{regola}

% ============================================================
% SEZIONE: RISOLUZIONE DELLA SENSIBILE
% (Integrazione da armoniainrete.altervista.org/armonia-lezione-09)
% ============================================================
\section{La risoluzione della sensibile}
\label{sec:risoluzione_sensibile}

La \textbf{sensibile} --- il VII grado della scala --- è il grado più ``direzionato'' dell'intero sistema tonale. Essa dista solo un semitono dalla tonica e, in virtù di questa vicinanza, possiede una forza attrattiva che non può essere ignorata nella condotta delle voci.

\begin{regola}[I tre movimenti possibili della sensibile]
	\label{reg:tre_movimenti_sensibile}
	La sensibile può muoversi in tre modi:
	\begin{enumerate}
		\item \textbf{Sale di semitono alla tonica} (risoluzione primaria e obbligatoria nel soprano). Questo è il movimento più naturale e il più atteso dall'orecchio.
		\item \textbf{Scende di terza alla dominante}, a condizione che si trovi in una \textbf{parte interna} dell'accordo (contralto o tenore, mai soprano) e che il basso proceda in \textbf{moto contrario} rispetto a essa. Questo espediente permette di ottenere un accordo di tonica completo, evitando di triplicare la fondamentale.
		\item \textbf{Scende di semitono cromatico} (abbassamento della sensibile): questo movimento è possibile solo in presenza di un \textbf{cambio di tonalità o di una tonicizzazione}. La sensibile abbassata non è più sensibile della tonalità d'impianto, ma diventa un altro grado funzionale nella nuova tonalità locale.
	\end{enumerate}
	Ogni altro movimento della sensibile è da evitare, a meno che essa si trovi in una voce intermedia e il movimento sia giustificato dal contesto armonico.
\end{regola}

\begin{esempio}[Le tre risoluzioni della sensibile in Do Maggiore]
	\label{ex:risoluzioni_sensibile}

	\textbf{A. Risoluzione primaria (sensibile sale alla tonica):}
	\musica{es_cap04_01}

	\vspace{0.5cm}
	\textbf{B. Sensibile che scende di terza (voce interna, moto contrario con il basso):}
	\musica{es_cap04_02}

	In questo caso il Si (sensibile) si trova al contralto e scende di terza al Sol, mentre il basso sale da Sol a Do. Il moto contrario e la posizione interna della sensibile rendono il collegamento corretto.
\end{esempio}

\begin{errore}[Sensibile non risolta nel soprano]
	\label{err:sensibile_non_risolta}
	Se la sensibile si trova al \textbf{soprano}, essa \textbf{deve} salire di semitono alla tonica. Non è mai ammesso che il soprano scenda di terza, anche se tecnicamente la nota di arrivo appartiene all'accordo. Questo errore produce un soprano non cantabile e armonicamente incoerente.
\end{errore}

% ============================================================
% SEZIONE: ERRORI ARMONICI
% ============================================================
\section{Errori armonici}
\label{sec:errori_armonici}

Gli errori armonici si producono nel \textbf{collegamento tra due voci} e riguardano la qualità degli intervalli che le voci formano tra loro in successione. Si dividono in errori gravi ed errori di media gravità.

\subsection{Errori armonici gravi: quinte e ottave parallele}

\begin{errore}[Quinte e Ottave parallele]
	\label{err:quinte_ottave_parallele}
	È severamente vietato muovere due voci per moto parallelo mantenendo tra loro un intervallo di \textbf{quinta giusta} o di \textbf{ottava giusta}. Questo errore distrugge l'illusione di indipendenza delle parti, facendole percepire come un'unica linea melodica raddoppiata.
	\musica{es_cap04_03}
\end{errore}

Due ottave o due quinte sono ammesse se sono \textbf{separate da un'armonia diversa} e non sono quindi in successione diretta.

\begin{regola}[Quinte parallele Giusta – Diminuita]
	\label{reg:quinte_g_dim}
	È ammesso il movimento di quinte parallele quando la \textbf{prima quinta è giusta} e la \textbf{seconda è diminuita}. Questo movimento è accettato perché la quinta diminuita, essendo dissonante, non viene percepita come una vera quinta parallela.
	\musica{es_cap04_04}
\end{regola}

\subsection{Errori armonici meno gravi: quinte e ottave nascoste}

\begin{errore}[Quinte e ottave nascoste (o dirette)]
	\label{err:quinte_nascoste}
	Si deve evitare di raggiungere un'ottava o una quinta giusta tra le parti estreme (Soprano e Basso) muovendo entrambe le voci nella stessa direzione per salto.
\end{errore}

\begin{regola}[Eccezioni alle quinte e ottave nascoste]
	\label{reg:eccezioni_nascoste}
	Le quinte e ottave nascoste tra parti estreme sono tollerate nei seguenti casi:
	\begin{itemize}
		\item Per le \textbf{ottave nascoste}: il basso sale di almeno una quarta giusta e il soprano sale di semitono diatonico (tipicamente la sensibile che sale alla tonica).
		\item Per le \textbf{quinte nascoste}: una delle voci si muove per grado congiunto, oppure il basso si muove verso la tonica o la dominante e la voce superiore procede per grado.
	\end{itemize}
\end{regola}

\subsection{Incrocio delle parti}

\begin{errore}[Incrocio delle voci]
	\label{err:incrocio}
	È proibito che una voce più grave canti una nota più acuta della voce adiacente superiore. Il Tenore non può cantare una nota più acuta del Contralto, né il Contralto una più acuta del Soprano.
\end{errore}

% ============================================================
% SEZIONE: SINCOPE ARMONICA
% (Integrazione da armoniainrete lezione-08)
% ============================================================
\section{La sincope armonica}
\label{sec:sincope_armonica}

Un caso particolare di condotta delle voci è la \textbf{sincope armonica}: una voce tiene una nota che nell'accordo successivo risulta dissonante rispetto alle altre voci. Questo fenomeno è strettamente legato all'uso dei ritardi (Capitolo \ref{chap:ritardi}) ed è uno dei pilastri espressivi dello stile corale.

\begin{definizione}[Sincope armonica]
	\label{def:sincope_armonica}
	La \textbf{sincope armonica} si verifica quando una nota \textit{consonante} rispetto all'accordo corrente viene \textit{tenuta} nel momento in cui il basso si muove verso un nuovo accordo, diventando così \textit{dissonante} rispetto alla nuova armonia.
\end{definizione}

\begin{regola}[Trattamento della sincope armonica]
	\label{reg:sincope_armonica}
	La nota tenuta per sincope deve:
	\begin{enumerate}
		\item Essere preparata come consonanza nell'accordo precedente.
		\item Cadere sul tempo forte (o sulla parte forte del tempo), rafforzando la tensione.
		\item Risolvere per grado congiunto \textbf{discendente} verso la nota reale dell'accordo (in questo modo si comporta come un ritardo).
	\end{enumerate}
\end{regola}

% ============================================================
% SEZIONE: RADDOPPI
% ============================================================
\section{I raddoppi}
\label{sec:raddoppi}

Poiché scriviamo a quattro parti ma gli accordi hanno tipicamente tre note, una di esse deve essere raddoppiata (o triplicata nella risoluzione della V7 completa).

\begin{regola}[Principio generale dei raddoppi]
	\label{reg:raddoppi}
	\begin{itemize}
		\item \textbf{Raddoppiare di norma la fondamentale}, specialmente negli stati fondamentali.
		\item \textbf{Mai} raddoppiare la \textbf{sensibile} (VII grado), perché porta inevitabilmente a ottave parallele nella risoluzione (entrambe le sensibili devono salire alla tonica).
		\item \textbf{Mai} raddoppiare la \textbf{settima} degli accordi di settima.
		\item Nei rivolti, è spesso utile raddoppiare la \textbf{terza} (specialmente nel collegamento II7-V per evitare quinte parallele, cfr. \ref{reg:ii7_raddoppio}).
		\item Il III grado maggiore è un grado caratteristico e va raddoppiato con parsimonia.
	\end{itemize}
\end{regola}

% ============================================================
% SEZIONE: CHECKLIST REVISIONE
% ============================================================
\section{Lista di controllo: revisione dell'elaborato}

\begin{regola}[Checklist errori]
\label{reg:checklist_errori}
Una volta completata l'armonizzazione, sottoporre l'esercizio a questa revisione sistematica.

\textbf{Blocco 1 — Errori verticali (accordo per accordo):}
\begin{enumerate}
	\item Ogni accordo contiene le note corrette per la sua funzione?
	\item Il raddoppio è corretto? (Raddoppiare la fondamentale di default; \textbf{mai} raddoppiare sensibile o settima).
	\item Le voci rispettano la loro estensione? (\ref{def:estensione_voci}).
	\item Le distanze tra voci adiacenti sono corrette? (\ref{reg:distanza_voci}).
\end{enumerate}

\textbf{Blocco 2 — Errori orizzontali (voce per voce):}
\begin{enumerate}
	\item Ci sono quinte parallele tra qualsiasi coppia di voci?
	\item Ci sono ottave parallele tra qualsiasi coppia di voci?
	\item Ci sono quinte o ottave nascoste tra le voci esterne?
	\item La sensibile nel soprano risolve salendo alla tonica? (\ref{reg:tre_movimenti_sensibile}).
	\item La sensibile in voce interna, se scende di terza, lo fa in moto contrario con il basso?
	\item La settima risolve scendendo di grado?
	\item Ci sono salti non giustificati (sesta, settima) in una voce?
	\item Ci sono movimenti di seconda aumentata (specialmente nel minore)?
	\item Ci sono false relazioni tra voci diverse? (\ref{def:falsa_relazione}).
	\item Ci sono tritoni armonici tra voci non adiacenti in successione?
\end{enumerate}

\textbf{Blocco 3 — Errori strutturali (visione d'insieme):}
\begin{enumerate}
	\item Le cadenze sono posizionate nei punti metricamente corretti?
	\item Le frasi hanno una struttura riconoscibile? (\ref{chap:struttura_frase}).
	\item Il basso dato e l'eventuale numerica sono stati rispettati?
\end{enumerate}
\end{regola}

Si suggerisce di rileggere ogni \textbf{coppia} di voci isolata (S+B, A+T, S+A, T+B ecc.) anziché leggere tutto l'elaborato insieme: gli errori di parallelismo si individuano molto più facilmente con questa tecnica.
